\Askhsh{13699}{4}
\begin{ylh}{1}
    \item 3.11. Ανισοτικές σχέσεις πλευρών και γωνιών
    \item 3.14. Σχετικές θέσεις ευθείας και κύκλου
    \item 4.2. Τέμνουσα δύο ευθειών - Ευκλείδειο αίτημα
    \item 5.3. Ορθογώνιο.
\end{ylh}
% ===================================================================
%  Αρχείο: 13699.tex (Fragment)
%  Αυτόματη μετατροπή μέσω doc2latex.py
% ===================================================================

ΘΕΜΑ 4

Δίνονται δυο κύκλοι (Κ, ρ\textsubscript{1}) και (Λ, ρ\textsubscript{2}) που εφάπτονται εξωτερικά σε σημείο Α. Έστω ότι μια ευθεία (ε) εφάπτεται εξωτερικά στους δυο κύκλους σε σημεία τους Β και Γ αντίστοιχα και ότι η εσωτερική εφαπτομένη (ζ) των κύκλων στο σημείο επαφής τους Α τέμνει την ευθεία (ε) σε σημείο Μ.

\begin{alist}
\item Να αποδείξετε ότι:

.}
\end{alist}
\begin{alist}
\item
  οι ευθείες ΚΒ και ΛΜ τέμνονται σε σημείο, έστω Δ. \monades{10}
\item
  το τρίγωνο ΔΚΛ είναι ισοσκελές. \monades{10}

\item Με ποια σχέση πρέπει να συνδέονται οι ακτίνες ρ\textsubscript{1} και ρ\textsubscript{2} των δύο κύκλων ώστε το ισοσκελές τρίγωνο ΔΚΛ να είναι ορθογώνιο; Να αιτιολογήσετε την απάντησή σας.

\monades{5}
\end{alist}
